\scanpage{56}
ce qui fait qu'il suffit de définir l'un des deux isom. (34) -- et en vertu de la définition de $\Pi_{03}$ dans (32), il est naturel de définir plutôt l'iso
\[
\deleted{(\unsure{35})}\qquad\psi_r=\psi_{\deleted{\unsure{\omega}}\varepsilon}^{i,\omega}:\Pi_{0,3}\xrightarrow{\sim}\Pi_\varepsilon
\]
\margin{$r=(i,\omega,\underset{\deleted{\unsure{\varepsilon}}}{\overbrace{\varepsilon,\omega}^{\bar\omega}})$}
par transport de structure comme
\[
(\unsure{36|35})\quad\left\{\begin{array}{l}
\psi_{\varepsilon\deleted{\unsure{\omega}}}^{i,\omega}=\deleted{\Pi_1(\unsure{u_{\varepsilon\omega}^{i,\omega}|u_\varepsilon^{i,\omega}})}\Pi_1(u_{i,\omega,\varepsilon};P^+)\\
\text{où }\inserted{u_r=u_{i,\omega,\varepsilon}} u_{\varepsilon\deleted{\unsure{\omega}}}^{i,\omega}:(U_{0,3}(\mathbb C)\unsure{,|;}P^{\deleted{\unsure{+}}}_{+})\longrightarrow(\Sigma^*,P_\omega).\\
\text{est l'unique iso conforme qui transforme}\\
\text{le repère \deleted{\unsure{base}} }\text{\unsure{standard|standart}}\ \inserted{\unsure{[|\{}\infty,\omega_{01\infty},+1\unsure{]|\}}}\ \text{\deleted{\unsure{analogue}} dans le}
\end{array}\right.
\]
repère $\deleted{\unsure{(\omega^2,i)|(\omega i,i)}}$ $(i,\omega,\varepsilon)$
\margin{Au couple $(i,\omega)$ on associe le repère (au sens habituel) $r=(\omega i,\omega\prime i,i)$ de $J$ i.e. \deleted{\unsure{correspondant}}\[
\begin{gathered}
0\longmapsto\omega i\\
1\longmapsto\omega^2i=\omega\prime i\\
\infty\longmapsto i
\end{gathered}
\]
Ainsi, on \unsure{a|aura}, si $r=\deleted{\unsure{(\omega i,\omega\prime i,i)}}\inserted{(i,\omega)}$\[
\rho_\omega(r)=r\circ\rho=(\omega i,\omega)
\]}

Pour un $P_\omega$ fixé, il y a six choix pour les $\inserted{\psi_r=}\ \psi_{\varepsilon\deleted{\unsure{\omega}}}^{i,\omega}:\Pi_{03}\longrightarrow\Pi_\varepsilon$ -- et pour un $Q_i$ fixé, il y a quatre choix pour $\omega,\varepsilon$. Cela fait \inserted{trois resp.} deux choix si on suppose que l'on a choisi une valeur de $\uline{\bar\omega}$ sur $\Sigma$ et réalisé donc $\uline{\varepsilon}\simeq\uline{\omega}$, et si on se borne aux $\varepsilon$ qui correspondent \unsure{à|:} $\omega\ \unsure{\Rightarrow|)}$ -- il faut donc voir comment passer de l'un \inserted{de ces choix} à l'autre.
\margin{55}
